\documentclass[12pt]{article}
\usepackage{amsmath, amssymb, amsthm}
\usepackage{geometry}
\usepackage{xcolor}
\usepackage{graphicx}
\usepackage{nicematrix}
\usepackage{multicol}
\usepackage{booktabs}
\usepackage{tikz}
\usepackage{sectsty}
\usepackage{cancel}
\geometry{margin=1in}

\title{Engineering Probability \\ Homework 8}
\author{Timothy Pollard, pollat@rpi.edu}
\date{October 25th 2025}

\sectionfont{\large\bfseries\color{black}} % makes them blue and large
\newenvironment{bluetext}
    {\par\begingroup\color{blue}\parindent=0pt\ignorespaces}
    {\par\endgroup\noindent\ignorespacesafterend}
\begin{document}
\maketitle
\hrule
\section*{(1) Joint PMF of Two Discrete Random Variables.}

Random variables $X$ and $Y$ have the joint PMF

\[
P_{X,Y}(x,y) =
\begin{cases}
c\,|x+y|, & x = -2, 0, 2;~ y = -1, 0, 1, \\
0, & \text{otherwise.}
\end{cases}
\]

\begin{enumerate}
    \item[(a)] Find the value of the constant $c$.
    \begin{bluetext}
        \[\sum_{S_{X,Y}}^{}P_{X,Y}(x,y)=1\]
    \[P(-2,-1)+ P(-2,0)+ P(-2,1)+ 
     P(0,-1)+ P(0,0)+ P(0,1) +
     P(2,-1)+ P(2,0)+ P(2,1)=1 \]
     \[c|(-2-1)|+ c|(-2+0)|+ c|(-2+1)|+ 
     c|(0-1)|+ c|(0+0)|+ c|(0+1)| +
     c|(2-1)|+ c|(2+0)|+ c|(2+1)|=1 \]
     \[c\left(|(-3)|+ |(-2)|+ |(-1)|+ 
     |(-1)|+ |(0)|+ |(1)| +
     |(1)|+ |(2)|+ |(3)|\right)=1 \]
     \[c\left(3+ 2+ 1+ 
     1+ 0+ 1 +
     1+ 2+ 3\right)=1 \]
     \[c\left(14\right)=1, c=\frac{1}{14} \]
    \end{bluetext}
    \item[(b)] Compute $P[Y < X]$.
    \begin{bluetext}
        \[
        \begin{aligned}
            P[Y<X]&=P(0,-1)+P(2,-1)+P(2,0)+P(2,1)\\
            P[Y<X]&= \frac{1}{14}\left(|(0-1)|+|(2-1)|+|(2+0)|+|(2+1)|\right)\\
            P[Y<X]&= \frac{1}{14}\left(1+1+2+3\right)\\
            P[Y<X]&=\frac{7}{14}\\
        \end{aligned}
        \]
    \end{bluetext}
    \subsection*{(1) Continued.}
    \item[(c)] Compute $P[Y > X]$.
    \begin{bluetext}
        \[
        \begin{aligned}
            P[Y>X]&=P(-2,-1)+P(-2,0)+P(-2,1)+P(0,1)\\
            P[Y>X]&= \frac{1}{14}\left(|(-2-1)|+|(-2+0)|+|(-2+1)|+|(0+1)|\right)\\
            P[Y>X]&= \frac{1}{14}\left(3+2+1+1\right)\\
            P[Y>X]&=\frac{7}{14}\\
        \end{aligned}
        \]
    \end{bluetext}
    \item[(d)] Compute $P[Y = X]$.
    \begin{bluetext}
        \[P[Y=X]=P(0,0)=\frac{1}{14}\left(|(0+0)|\right)=0\]
    \end{bluetext}
    \item[(e)] Compute $P[X < 1]$.
    \begin{bluetext}
        \[
        \begin{aligned}
            P[X<1]&=P[X=-2]+P[X=0]\\
            &=P(-2,-1)+P(-2,0)+P(-2,1)+P(0,-1)+P(0,0)+P(0,1)\\
            &=\frac{1}{14}\left(|(-2-1)|+|(-2+0)|+|(-2+1)|+|(0-1)|+|(0+0)|+|(0+1)|\right)\\
            &=\frac{1}{14}\left(3+2+1+1+0+1\right)\\
            &=\frac{8}{14}=\frac{4}{7}
        \end{aligned}
        \]
    \end{bluetext}    
\end{enumerate}
\newpage
\section*{(2) Testing Two Circuits.}

Two integrated circuits are tested. In each test, the probability of rejecting the circuit is $p$. Let $X$ be the number of rejects in the first test and $Y$ the number of rejects in the second test.

\begin{enumerate}
    \item[(a)] Find the joint PMF $P_{X,Y}(x, y)$.
    \begin{bluetext}
        Possible outcomes for (X,Y): (0,0), (0,1), (1,0), (1,1)
        \[P_X(0)=1-p, P_X(1)=p\]
        \[P_Y(0)=1-p, P_Y(1)=p\]
        \[
        P_{X,Y}(x,y)=
        \begin{cases}
        p^{x+y}(1-p)^{2-(x+y)} & x=0,1;~ y=0,1, \\
        0, & \text{otherwise.}
        \end{cases}
        \]
    \end{bluetext}
\end{enumerate}

\section*{(3) Rejects and Acceptables.}

Each test of an integrated circuit produces an acceptable circuit with probability $p$, independent of other tests. Let $N$ denote the number of circuits rejected in testing $r$ circuits. Let $X$ denote the number of acceptable circuits that appear before the first reject is found.

\begin{enumerate}
    \item[(a)] Find the joint PMF $P_{N,X}(n, x)$.
    \begin{bluetext}
        \[x=r, n=0\]
        \[P_{N,X}(0,r)=p^r\]
        \[P_{N,X}(n,x)=
        \begin{cases}
        p^r, & n=0,\ x=r,\\
        \displaystyle \binom{r-x-1}{\,n-1\,}\,p^{\,r-n}\,(1-p)^{\,n}, & x=0,1 \cdots r-1, n=1\cdots r-x \\
        0, & \text{otherwise.}
        \end{cases}
        \]
    \end{bluetext}
\end{enumerate}
\newpage
\section*{(4) Marginals and Moments.}

Given the random variables $X$ and $Y$ in Problem 1:

\begin{enumerate}
    \item[(a)] Find the marginal PMFs $P_X(x)$ and $P_Y(y)$.
    \begin{bluetext}
        \[
P_{X,Y}(x,y) =
\begin{array}{c|ccc}
x \backslash y & -1 & 0 & 1 \\ \hline
-2 & \tfrac{3}{14} & \tfrac{2}{14} & \tfrac{1}{14} \\[4pt]
0  & \tfrac{1}{14} & 0             & \tfrac{1}{14} \\[4pt]
2  & \tfrac{1}{14} & \tfrac{2}{14} & \tfrac{3}{14}
\end{array}
\]
        \[P_X(x)= 
        \begin{cases}
        \frac{6}{14}, & x=-2\\
        \frac{2}{14}, & x=0\\
        \frac{6}{14}, & x=2  \\
        0, & \text{otherwise.}
        \end{cases}
        \qquad
        P_Y(y)=
        \begin{cases}
        \frac{5}{14}, & y=-1\\
        \frac{4}{14}, & y=0\\
        \frac{5}{14}, & y=1  \\
        0, & \text{otherwise.}
        \end{cases}\]
    \end{bluetext}
    \item[(b)] Compute the expected values $E[X]$ and $E[Y]$.
    \begin{bluetext}
        \[
            E[X]=(-2)\left(\frac{6}{14}\right)+(0)\left(\frac{2}{14}\right)+(2)\left(\frac{6}{14}\right)=0
        \]
        \[
            E[Y]=(-1)\left(\frac{5}{14}\right)+(0)\left(\frac{4}{14}\right)+(1)\left(\frac{5}{14}\right)=0
        \]
    \end{bluetext}

    \item[(c)] Compute the standard deviations $\sigma_X$ and $\sigma_Y$.
    \begin{bluetext}
        \[\sigma_X=\sqrt{\text{Var}[X]}=\sqrt{E[X^2]-\left(E[X]\right)^2}\qquad\sigma_Y=\sqrt{\text{Var}[Y]}=\sqrt{E[Y^2]-\left(E[Y]\right)^2}\]
        \[
            E[X^2]=(-2)^2\left(\frac{6}{14}\right)+(0)^2\left(\frac{2}{14}\right)+(2)^2\left(\frac{6}{14}\right)
            =\frac{48}{14}=\frac{24}{7}
        \]
        \[
            E[Y^2]=(-1)^2\left(\frac{5}{14}\right)+(0)^2\left(\frac{4}{14}\right)+(1)^2\left(\frac{5}{14}\right)
            =\frac{10}{14}=\frac{5}{7}
        \]  
         \[\sigma_X=\sqrt{\left(\frac{24}{7}\right)-\left(0\right)^2}\qquad\sigma_Y=\sqrt{\left(\frac{5}{7}\right)-\left(0\right)^2}\]
         \[\sigma_X=\sqrt{\frac{24}{7}}\qquad\sigma_Y=\sqrt{\frac{5}{7}}\]
         \[\sigma_X\approx1.852\qquad\sigma_Y\approx0.845\]
    \end{bluetext}
\end{enumerate}

\section*{(5) Function of Two Variables.}

For the random variables $X$ and $Y$ in Problem 1, define $W = X + 2Y$.

\begin{enumerate}
    \item[(a)] Find the probability mass function $P_W(w)$.
    \begin{bluetext}
        Possible values of $w$:
        \[
\begin{array}{c|c|c|c}
x & y & P(x,y) & w \\ \hline
-2 & -1 & \frac{3}{14} & -4 \\[4pt]
-2 & 0 & \frac{2}{14} & -2 \\[4pt]
-2 & 1 & \frac{1}{14} & 0 \\[4pt]
0  & -1 & \frac{1}{14} & -2 \\[4pt]
0  & 0 & \frac{0}{14} & 0 \\[4pt]
0  & 1 & \frac{1}{14} & 2 \\[4pt]
2  & -1 & \frac{1}{14} & 0 \\[4pt]
2  & 0 & \frac{2}{14} & 2 \\[4pt]
2  & 1 & \frac{3}{14} & 4 \\
\end{array}
\qquad P_W(w)= \begin{cases}
        \frac{3}{14}, & x=-4\\
        \frac{3}{14}, & x=-2\\
        \frac{2}{14}, & x=0  \\
        \frac{3}{14}, & x=2  \\
        \frac{3}{14}, & x=4  \\
        0, & \text{otherwise.}
        \end{cases}\]
    \end{bluetext}
    \item[(b)] Compute $E[W]$.\
    \begin{bluetext}
        \[
            E[W]=(-4)\left(\frac{3}{14}\right)+(-2)\left(\frac{3}{14}\right)+(0)\left(\frac{2}{14}\right)+(2)\left(\frac{3}{14}\right)+(4)\left(\frac{3}{14}\right)=0
        \]
    \end{bluetext}
    \item[(c)] Compute $P[W > 0]$.
    \begin{bluetext}
        \[
        \begin{aligned}
            P[W>0]&=P_W(2)+P_W(4)\\
            P[W>0]&=\frac{3}{14}+\frac{3}{14}\\
            P[W>0]&=\frac{6}{14}\\
            P[W>0]&=\frac{3}{7}\\
        \end{aligned}
        \]
    \end{bluetext}
\end{enumerate}
\newpage
\section*{(6) Joint Dependence and Correlation.}

For the random variables $X$ and $Y$ in Problem 1:
\begin{bluetext}
    \[
P_{X,Y}(x,y) =
\begin{cases}
\frac{1}{14}\,|x+y|, & x = -2, 0, 2;~ y = -1, 0, 1, \\
0, & \text{otherwise.}
\end{cases}
\]
\end{bluetext}
\begin{enumerate}
    \item[(a)] Compute the expected value of $W = 2XY$.
        \begin{bluetext}
        \[
        \begin{array}{c|c|c|c}
x & y & P(x,y) & w \\ \hline
-2 & -1 & \frac{3}{14} & 4 \\[4pt]
-2 & 0 & \frac{2}{14} & 0 \\[4pt]
-2 & 1 & \frac{1}{14} & -4 \\[4pt]
0  & -1 & \frac{1}{14} & 0 \\[4pt]
0  & 0 & \frac{0}{14} & 0 \\[4pt]
0  & 1 & \frac{1}{14} & 0 \\[4pt]
2  & -1 & \frac{1}{14} & -4 \\[4pt]
2  & 0 & \frac{2}{14} & 0 \\[4pt]
2  & 1 & \frac{3}{14} & 4 \\
\end{array}\qquad
        P_{W}(w)=
        \begin{cases}
        \frac{6}{14}, & w=-4\\
        \frac{6}{14}, & w=0\\
        \frac{2}{14}, & w=4  \\
        0, & \text{otherwise.}
        \end{cases}
        \]
            \[
            E[W]=\sum_{S_{W}}wP_{W}(w)=\left(-4\right)\frac{6}{14}+\left(0\right)\frac{6}{14}+\left(4\right)\frac{2}{14}=\frac{-16}{14}=\frac{-8}{7}
            \]
        \end{bluetext}
    \item[(b)] Compute the correlation $E[XY]$.
    \begin{bluetext}
        \[
                E[XY]=\frac{1}{2}E[2XY]=\frac{1}{2}\left(\frac{-8}{7}\right)=\frac{-4}{7}
        \]
    \end{bluetext}
    \item[(c)] Compute the covariance $\mathrm{Cov}[X, Y]$.
        \begin{bluetext}
            \[
            \mathrm{Cov}[X,Y]=E[XY]-E[X]E[Y]= \frac{-4}{7}-\left(0\right)\left(0\right)=\frac{-4}{7}
            \]
        \end{bluetext}
    \item[(d)] Compute the correlation coefficient $\rho_{X,Y}$.
        \begin{bluetext}
            \[
            \rho_{X,Y}=\frac{\mathrm{Cov}[X,Y]}{\sigma_X\sigma_Y}=\frac{\frac{-4}{7}}{\left(\sqrt{\frac{24}{7}}\right)\left(\sqrt{\frac{5}{7}}\right)}=\frac{-4}{\sqrt{120}}\approx-0.365
            \]
        \end{bluetext}
    \subsection*{(6) Continued.}
    \item[(e)] Compute the variance $\mathrm{Var}[|X + Y|]$.
        \begin{bluetext}
            \[
\begin{array}{c|c|c|c}
x & y & P(x,y) & z \\ \hline
-2 & -1 & \frac{3}{14} & 3 \\[4pt]
-2 & 0 & \frac{2}{14} & 2 \\[4pt]
-2 & 1 & \frac{1}{14} & 1 \\[4pt]
0  & -1 & \frac{1}{14} & 1 \\[4pt]
0  & 0 & \frac{0}{14} & 0 \\[4pt]
0  & 1 & \frac{1}{14} & 1 \\[4pt]
2  & -1 & \frac{1}{14} & 1 \\[4pt]
2  & 0 & \frac{2}{14} & 2 \\[4pt]
2  & 1 & \frac{3}{14} & 3 \\
\end{array}\qquad P_Z(z)=\begin{cases}
        \frac{4}{14}, & z=1\\
        \frac{4}{14}, & z=2  \\
        \frac{6}{14}, & z=3  \\
        0, & \text{otherwise.}
\end{cases}\]
\[E[Z]=E[|X+Y|]=1\cdot\frac{4}{14}+2\cdot\frac{4}{14}+3\cdot\frac{6}{14}=\frac{30}{14}\]
\[E[Z^2]=E[|X+Y|^2]=1^2\cdot\frac{4}{14}+2^2\cdot\frac{4}{14}+3^2\cdot\frac{6}{14}=\frac{74}{14}\]
            \[
            \mathrm{Var}[|X + Y|]=E[|X + Y|^2]-\left(E[|X + Y|]\right)^2=\frac{74}{14}-\left(\frac{30}{14}\right)^2=\frac{34}{49}\approx0.694
            \]
        \end{bluetext}
\end{enumerate}
\newpage
\section*{(7) Dependent Support Random Variables.}

Random variables $X$ and $Y$ have the joint PMF

\[
P_{X,Y}(x,y) =
\begin{cases}
\frac{1}{21}, & x = 0, 1, 2, 3, 4, 5;~ y = 0, 1, \ldots, x, \\
0, & \text{otherwise.}
\end{cases}
\]
\begin{enumerate}
    \item Find the marginal PMFs $P_X(x)$ and $P_Y(y)$.
    \begin{bluetext}
        all possible values of (x,y): Are equally likely. Simply a matter of counting number of x=0s, x=1s etc and same for y. 
        x can take any number of values of y up to x so its simply x plus the value itself or x+1.
        y can take any number of values of x so it is total number of x values (6) minus y.
        \[P_X(x)=\begin{cases}
        \frac{x+1}{21}, & x = 0, 1, 2, 3, 4, 5; \\
        0, & \text{otherwise.}
        \end{cases}\qquad P_Y(y)=\begin{cases}
        \frac{6-y}{21}, & y = 0, 1, 2, 3, 4, 5; \\
        0, & \text{otherwise.}
        \end{cases}\]
    \end{bluetext}
    \item Find the covariance $\mathrm{Cov}[X, Y]$.
    \begin{bluetext}
        \[
        \begin{aligned}
            E[XY]&=\sum_{S_{X,Y}}xyP_{X,Y}(x,y)\\
            E[XY]&=\sum_{0}^{5}\sum_{0}^{x}xy\left(\frac{1}{21}\right)=\frac{1}{21}\sum_{0}^{5}x\left(\frac{x(x+1)}{2}\right)=\frac{1}{42}\sum_{0}^{5}x^2+x^3\\
            E[XY]&=\frac{1}{42}\left((0)^2+(0)^3+(1)^2+(1)^3+(2)^2+(2)^3+(3)^2+(3)^3+(4)^2+(4)^3+(5)^2+(5)^3\right)\\
            E[XY]&=\frac{1}{42}\left(1+1+4+8+9+27+16+64+25+125\right)=\frac{280}{42}\\
        \end{aligned}
        \]
\[
\begin{aligned}
E[X]&=\sum_{S_X} xP_X(x)\\
E[X]&=\sum_{x=0}^{5} x\left(\frac{x+1}{21}\right)
    =\frac{1}{21}\sum_{x=0}^{5}(x^2+x)\\
E[X]&=\frac{70}{21}=\frac{10}{3}
\end{aligned}
\qquad
\begin{aligned}
E[Y]&=\sum_{S_Y} yP_Y(y)\\
E[Y]&=\sum_{y=0}^{5} y\left(\frac{6-y}{21}\right)
    =\frac{1}{21}\sum_{y=0}^{5}(6y-y^2)\\
E[Y]&=\frac{35}{21}=\frac{5}{3}
\end{aligned}
\]

\[
\mathrm{Cov}[X,Y]=\left(\frac{280}{42}\right)-\left(\frac{10}{3}\right)\left(\frac{5}{3}\right)=\frac{10}{9}\approx1.11\\
\]        
        
    \end{bluetext}
\end{enumerate}

\end{document}

